% K17: low-perfection-excess simpliciality and partial target-cell pruning.
% Author /root/vandiver_k; GPT-6 Astra Pro; effort not exposed.
% 2026-11-21 UTC. This section proves no full K12 vanishing.

\subsection{Simpliciality and partial orientation pruning of the remaining cells}
\label{subsec:k12-orientation}

For a well-rounded lattice $L$ of rank $N$, write $p(v)=vv^t$, let
$r$ be the dimension of the real span of its minimal-vector quadratic
forms, and choose one representative of each minimal pair.
The corresponding projective cell has dimension $r-1$.
The maximal index $\iota(L)$ is the maximum of
$[L:\sum_{j=1}^N\mathbb Z v_j]$ over independent minimal vectors.
The following deduction uses the known theorem that every perfect cone
of dimension at most nine is basic and therefore simplicial
\cite[Theorem~1.1(i), pp.~644--646]{DutourHulekSchurmann2015}.
Here the cited dimension is the dimension of the cone's linear span.

\begin{proposition}[K17-K12-CIRCUIT-SIMPLE]
\label{prop:k12-low-excess-simplicial}
If $r-N\leq4$, the minimal-vector quadratic forms are linearly
independent. In particular, every projective degree-twelve cell in
rank nine or ten has exactly thirteen minimal-vector pairs.
\end{proposition}

\begin{proof}
Put $d=r-N\geq0$. If the forms are dependent, choose an
inclusion-minimal dependent family $p(v_1),\ldots,p(v_k)$.
It has rank $k-1$, and its unique relation up to scaling has every
coefficient nonzero. Set $W=\operatorname{span}_{\mathbb R}
\{v_1,\ldots,v_k\}$ and $m=\dim W$. Separate the positive and negative
coefficients to obtain an equality of positive semidefinite sums.
The image of either sum is the span of its supporting vectors, so
both supporting sets span $W$. Each contains at least $m$ vectors;
therefore $k\geq2m$.

Let $\sigma_W$ be generated by all minimal-vector rays lying in $W$,
and let $r_W$ be its dimension. This is a face: if $P$ is orthogonal
projection onto $W^\perp$, the functional $q\mapsto\operatorname{tr}(Pq)$
is nonnegative on every generator and vanishes precisely on those
in $W$. The space $W$ is rational. Equivalently, $\sigma_W$ is the
perfect cone of $L\cap W$ with its restricted form; its minimum
equals that of $L$. The chosen circuit gives $r_W\geq k-1$.

Choose $N-m$ minimal vectors whose images form a basis of
$(L\otimes\mathbb R)/W$. Their projected quadratic forms are linearly
independent, since in suitable quotient coordinates they are the
diagonal elementary matrices. Projection annihilates the span of
$\sigma_W$, so
\[
 r\geq r_W+N-m\geq k-1+N-m\geq N+m-1.
\]
Thus $m\leq d+1$, and also
\[
 r_W\leq r-N+m=d+m\leq2d+1\leq9.
\]
The cited theorem makes $\sigma_W$ simplicial. Distinct minimal pairs
give distinct extremal rays: a positive sum of rank-one positive
semidefinite matrices has rank one only if every summand has the
same image. The circuit is therefore a dependence among distinct
extremal rays of $\sigma_W$, a contradiction.
\end{proof}

\begin{proposition}[K17-K12-ORIENT-LOWINDEX]
\label{prop:k12-low-index-orientation}
A projective degree-twelve cell has an orientation-reversing integral
stabilizer element in either of the following cases:
\begin{enumerate}
\item its lattice has rank nine and maximal index one;
\item its lattice has rank ten and maximal index at most two.
\end{enumerate}
Consequently these cells contribute no generators to the localized
Voronoi chain groups over $R=\mathbb Z[1/13!]$.
\end{proposition}

\begin{proof}
By Proposition~\ref{prop:k12-low-excess-simplicial}, there are thirteen
minimal pairs, and their quadratic forms are independent. First suppose
the maximal index is one. In a basis of minimal vectors write their
column matrix as $V=[I_N\ A]$. Every maximal minor has value $0,1$,
or $-1$, so $V$ is totally unimodular. The full integral relation
lattice is the row lattice of
\[
 B=[-A^t\ I_c],\qquad c=13-N,
\]
which is also totally unimodular. Heller's bound gives at most
$c(c+1)/2$ nonzero columns distinct up to sign
\cite[Theorem~4.2, pp.~1358--1360]{Heller1957}.
For $c=4$ or $3$ this bound is ten or six, respectively.
Thirteen columns therefore force either two zero columns or two
nonzero columns equal up to sign. After a sign change of a minimal
vector if necessary, their transposition fixes every row of $B$.
It descends to an integral involution of $L$, exchanges two minimal
rays, and fixes the others. Its determinant on their independent
quadratic span is $-1$.

It remains to treat rank ten with maximal index two. Let $L_0$ be
the lattice generated by all minima. If $L_0=L$, the gcd of the
maximal minors of $V$ is one. Since their nonzero values have
absolute value one or two, one minor has absolute value one.
Again use its columns as a basis and write $V=[I_{10}\ A]$.
The rank-three relation matrix $B=[-A^t\ I_3]$ has maximal minors
of absolute value at most two, by complementary minors. The exact
bimodular bound is
\[
 c(2,3)=\frac{3^2+3}{2}+3=9
\]
\cite[Theorem~1, p.~4]{LeePaatStallknechtXu2022}.
This counts nonzero columns distinct up to sign and permits
nonprimitive columns. Thirteen columns again force a signed-equal
pair or two zeros, producing the preceding integral ray transposition.

Suppose instead that $L_0\ne L$. Its index divides every independent
ten-vector index, hence $[L:L_0]=2$, and each independent ten minima
is a basis of $L_0$. Relative to one such basis, $V=[I_{10}\ A]$
and $B=[-A^t\ I_3]$ are totally unimodular. There are at most six
nonzero signed-column classes of $B$ and at most one zero class.
Orient the representatives so columns in each nonzero class are equal.
Distinct nonzero classes have distinct nonzero reductions in
$\mathbb F_2^3$: equal supports without a common overall sign would
give a $2\times2$ minor of absolute value two.

Write $L=L_0+\mathbb Z(w/2)$, where $w$ gives a nonzero class in
$L_0/2L_0$, and choose its coordinate lift $a\in\mathbb F_2^{13}$.
Integral surjectivity of $V$ gives
\[
 \ker(\overline V:\mathbb F_2^{13}\longrightarrow L_0/2L_0)
 =\operatorname{row}(\overline B).
\]
A permutation $P$ within the signed-column classes fixes the integral
relation lattice and defines an automorphism of $L_0$. It extends to
$L$ exactly when
\[
 Pa-a\in\operatorname{row}(\overline B).
\]
If a class has at least three columns, two corresponding entries of
$a$ agree, and their transposition fixes $a$. The same applies to a
two-column class with equal entries. Either yields the required
orientation reversal.

In the remaining case every class has size at most two and the
entries of $a$ in each doubled class are opposite. Thirteen columns
then force exactly seven classes, six nonzero and one zero, with six
doubles and one singleton. Their six nonzero labels omit a unique
nonzero $b\in\mathbb F_2^3$. Choose a linear functional $\lambda$
with $\lambda(b)=1$ and, if the singleton has a nonzero label $a_0$,
also $\lambda(a_0)=0$. This is possible because $a_0\ne b$.
A nonzero functional is one on four vectors of $\mathbb F_2^3$;
removing $b$ leaves exactly three present labels of value one,
all doubled. Let $P$ swap these three pairs. Since their glue bits
are opposite,
\[
 Pa-a=\lambda\overline B.
\]
Thus $P$ preserves the glue and extends to an integral automorphism
of $L$. It acts on the quadratic-ray basis as three transpositions,
so its determinant is $-1$.

In all cases the sum of the rays is a fixed positive radial vector;
passing to the projective cell preserves the orientation sign.
An orientation-reversing stabilizer makes the cell generator equal
to its negative in orientation coinvariants. Since two is invertible
in $R$, that generator is zero.
\end{proof}

These are partial removals of the two target chain groups. The
remaining cells have maximal index at least two in rank nine and
at least three in rank ten. Neither the proposition nor simpliciality
alone proves that the remaining homology groups vanish.
